\documentclass[11pt]{article}

\usepackage[margin=1in]{geometry}
\usepackage{amsmath,amssymb,amsthm}
\usepackage{graphicx}
\usepackage[hidelinks]{hyperref}

\newtheorem{theorem}{Theorem}
\newtheorem{corollary}{Corollary}
\newtheorem{remark}{Remark}

\newcommand{\FBL}{\operatorname{FBL}}
\newcommand{\Id}{\operatorname{Id}}

\title{Complemented Subspaces And Order Continuity In Free Banach Lattices}
\author{Automatic Math Discovery Run \texttt{fa\_banach\_001}}
\date{June 13, 2026}

\begin{document}
\maketitle

\section*{Source And Target}

Youssef Azouzi and Wassim Dhifaoui, \emph{Order denseness in free Banach
lattices}, arXiv:2508.11648.

The source paper revisits a theorem of Oikhberg--Taylor--Tradacete--Troitsky
asserting that, if \(F\) is a closed subspace of \(E\), then the
inclusion-induced homomorphism
\[
  \overline\iota:\FBL^{(p)}[F]\longrightarrow \FBL^{(p)}[E]
\]
is order continuous.  The paper explains that the old proof used order
denseness of \(\operatorname{FVL}[F]\) in \(\FBL^{(p)}[F]\), which holds only
when \(F\) is finite-dimensional.  It proves the statement for \(p=\infty\)
and says that, for other values of \(p\), the question remains open.

\begin{center}
\includegraphics[width=\textwidth]{figures/open_problem_crop.png}
\end{center}

\section*{Result}

\begin{theorem}
\label{thm:main}
Let \(1\leq p\leq\infty\).  Let \(T:F\to E\) be a bounded operator between
Banach spaces.  Suppose that \(T\) has a bounded left inverse, i.e. there is
a bounded operator \(R:E\to F\) such that \(RT=\Id_F\).  Then the induced
lattice homomorphism
\[
  \overline T:\FBL^{(p)}[F]\longrightarrow \FBL^{(p)}[E],
  \qquad
  \overline T f=f\circ T^*,
\]
is order continuous.
\end{theorem}

\begin{corollary}
If \(F\) is a complemented subspace of \(E\), then the natural
inclusion-induced homomorphism
\[
  \FBL^{(p)}[F]\longrightarrow \FBL^{(p)}[E]
\]
is order continuous for every \(1\leq p\leq\infty\).
\end{corollary}

\section*{Proof}

We use the standard functorial formula for free \(p\)-convex Banach lattices:
if \(A:X\to Y\) is bounded, then
\[
  \overline A:\FBL^{(p)}[X]\to \FBL^{(p)}[Y],
  \qquad
  \overline A f=f\circ A^*.
\]

Let \((f_\alpha)\) be a decreasing net in \(\FBL^{(p)}[F]_+\) with infimum
\(0\).  To prove order continuity of \(\overline T\), suppose that
\[
  0\leq h\leq \overline T f_\alpha
  \qquad(\alpha)
\]
for some \(h\in\FBL^{(p)}[E]_+\).  We show that \(h=0\).

Assume otherwise.  Since \(\FBL^{(p)}[E]\) is represented as a lattice of
positively homogeneous functions on \(E^*\), there is \(x^*\in E^*\) such
that \(h(x^*)>0\).  Put
\[
  y^*=T^*x^*\in F^*.
\]
Then \(y^*\neq0\), because if \(y^*=0\), then
\[
  h(x^*)\leq (\overline T f_\alpha)(x^*)=f_\alpha(0)=0.
\]

We now build a bounded operator \(S:E\to E\) satisfying
\[
  ST=T,
  \qquad
  S^*R^*y^*=x^*.
\]
Set \(d=x^*-R^*y^*\).  For every \(u\in F\),
\[
  d(Tu)=x^*(Tu)-y^*(RTu)=y^*(u)-y^*(u)=0.
\]
Choose \(u_0\in F\) such that \(y^*(u_0)=1\), and define
\[
  S e=e+d(e)\,T u_0 \qquad(e\in E).
\]
Then \(S\) fixes \(T(F)\), hence \(ST=T\).  Moreover, for every \(e\in E\),
\[
\begin{aligned}
  (S^*R^*y^*)(e)
  &= (R^*y^*)(Se)  \\
  &= y^*(Re+d(e)RTu_0) \\
  &= (R^*y^*)(e)+d(e)y^*(u_0) \\
  &= x^*(e).
\end{aligned}
\]

Apply the induced positive lattice homomorphisms \(\overline S\) and
\(\overline R\).  Since \(ST=T\), we have
\(\overline S\,\overline T=\overline T\).  Since \(RT=\Id_F\), we have
\(\overline R\,\overline T=\Id\) on \(\FBL^{(p)}[F]\).  Therefore, for every
\(\alpha\),
\[
  0\leq \overline R\,\overline S h
  \leq \overline R\,\overline S\,\overline T f_\alpha
  = \overline R\,\overline T f_\alpha
  = f_\alpha.
\]
But
\[
  (\overline R\,\overline S h)(y^*)
  = h(S^*R^*y^*)
  = h(x^*)>0.
\]
Thus \(\overline R\,\overline S h\) is a nonzero positive element of
\(\FBL^{(p)}[F]\) lying below every \(f_\alpha\), contradicting
\(f_\alpha\downarrow0\).  Hence \(h=0\), proving that
\(\overline T f_\alpha\downarrow0\).

\section*{Scope And Limitations}

This proves the finite-\(p\) closed-subspace order-continuity statement under
the additional assumption that the subspace is complemented.  Equivalently,
it proves the operator version when \(T\) is an embedding with complemented
range.

The proof does not settle the source paper's remaining open case for
arbitrary closed subspaces.  The point of the complemented hypothesis is that
the left inverse \(R\) lets us pull a hypothetical positive lower bound back
to \(\FBL^{(p)}[F]\).

\section*{Novelty And Duplicate Check}

Local duplicate check: the run registry, solution index, active claims, and
attempt logs were checked.  No packaged result covers arXiv:2508.11648.  A
prior search log identified the finite-\(p\) order-continuity question as
plausible but did not solve or promote any subcase.

No external literature search was performed for this partial packet.  The
novelty claim is therefore limited: this should be read as a compact
complemented-subspace proof extracted within the run, not as a claim that the
subcase is absent from the literature.

\section*{Verifier Focus}

The main points to check are:
\begin{enumerate}
\item the orientation of the functorial identities
\(\overline S\,\overline T=\overline T\) and
\(\overline R\,\overline T=\Id\);
\item the construction of \(S\), especially \(ST=T\) and
\(S^*R^*y^*=x^*\);
\item the order-continuity criterion using positive lower bounds for a
decreasing net.
\end{enumerate}

\end{document}
